\section{Kommunikation}
    \subsection{MQTT}
    Um eine Kommunikation zu ermöglichen, wurde überlegt, 
    eine MQTT-Library zu verwenden. Hierfür wurde "PubSubClient" von 
    "knolleary" verwendet. Hierbei handelt es sich um eine 
    simple Library, die nur das Nötigste implementiert, um
    einen überschaubaren Overhead zu gewährleisten.
        \subsubsection{Nachteile}
        Da MQTT eine WLAN-Verbindung für jeden ESP32
        benötigt und somit die Mobilität des Controlpanels
        deutlich einschränkt ist wurde entschieden, dass 
        ESP-NOW eine bessere Alternative darstellt.
    
    \subsection{ESP-NOW}
    Für eine Implementierung von ESP-NOW wurden die Libraries
    \texttt{esp-now.h} und \texttt{esp\_wifi.h} von
    „Espressif“ verwendet. Diese sind offiziell unterstützt
    und bieten eine einfache Möglichkeit, um eine Kommunikation
    zwischen ESP32's zu ermöglichen.

        \subsubsection{Eingabesicherung}
        Daten werden in Form von Structs übertragen. Diese Structs
        sind in der \texttt{FaRiLib} Library definiert. Um sicherzustellen,
        dass keine falschen Daten übertragen werden, wird gecheckt, ob die
        Strukturgröße mit der erwarteten Strukturgröße übereinstimmt. \par
        Diese Maßnahmen minimieren das Risiko von Fehlern der Werte oder unbefugten 
        Speicherzugriffen. Da das Struct direkt auf den Speicher des ESP32
        kopiert wird, können falsche Datentypen unvorhersehbare Fehler
        verursachen und anderen Speicherplatz überschreiben. \\

        Dieser Check kann wie folgt implementiert werden:
        \begin{lstlisting}[style=cppCode]
    void OnDataRecv(
        const uint8_t * mac, 
        onst uint8_t *incomingData, 
        int len) 
    {
        if(sizeof(espnowSlave.msgReceived) != len) return;
        //. . . Rest of the code
    }
        \end{lstlisting}
    
        Mit dieser Zeile wird die Datenverarbeitung abgebrochen,
        wenn die Strukturgröße nicht 
        übereinstimmt und der Speicher wird nicht überschrieben.